% Part: first-order-logic
% Chapter: tableaux
% Section: proof-theoretic-notions

\documentclass[../../../include/open-logic-section]{subfiles}

\begin{document}

\iftag{FOL}
      {\olfileid{fol}{tab}{ptn}}
      {\olfileid{pl}{tab}{ptn}}

\olsection{સાબિતી-સૈદ્ધાંતિક ખ્યાલો}

\begin{editorial}
આ વિભાગ ટેબ્લો માટે સાબિત કરી શકાય તેવા હોવાના સંબંધ અને સુસંગતતાની
વ્યાખ્યાઓ એકત્ર કરે છે.
\end{editorial}

\begin{explain}
આપણે અનેક મહત્વના અર્થાનુસારી ખ્યાલો (પ્રામાણ્ય, ફલિતતા અને સંતોષ્યતા)
વ્યાખ્યાયિત કર્યા છે; એ જ રીતે હવે તેમને અનુરૂપ \emph{સાબિતી-સૈદ્ધાંતિક
ખ્યાલો} વ્યાખ્યાયિત કરીએ છીએ. આ ખ્યાલો !!{structure}sમાં
!!{sentence}sના સંતોષનો આશરો લઈને નહિ, પરંતુ અમુક બંધ !!{tableau}xના
અસ્તિત્વનો આશરો લઈને વ્યાખ્યાયિત થાય છે. આ બંને પ્રકારના ખ્યાલો મળી જાય
છે એ એક મહત્વની શોધ હતી. તેઓ મળે છે તે જ \emph{યથાર્થતા} અને
\emph{પૂર્ણતાના પ્રમેયો}નો વિષય છે.
\end{explain}


\begin{defn}[પ્રમેયો]
$\sFmla{\False}{!A}$ માટે બંધ !!{tableau} હોય તો !!{sentence}~$!A$
\emph{પ્રમેય} છે. $!A$ પ્રમેય હોય તો $\Proves !A$ અને ન હોય તો
$\Proves/ !A$ લખીએ છીએ.
\end{defn}

\begin{defn}[!!^{derivability}]
!!^a{sentence} $!A$ એ !!{sentence}sના ગણ~$\Gamma$માંથી
\emph{!!{derivable}} છે, $\Gamma \Proves !A$, જો અને માત્ર જો એવો સાન્ત
ગણ $\{!B_1, \dots, !B_n\} \subseteq \Gamma$ અને નીચેના ગણ માટે બંધ
!!{tableau} હોય:
\[
\{
  \sFmla{\False}{!A},
  \sFmla{\True}{!B_1}, \dots,
  \sFmla{\True}{!B_n}
  \}.
\]
$!A$ એ $\Gamma$માંથી !!{derivable} ન હોય તો $\Gamma \Proves/ !A$
લખીએ છીએ.
\end{defn}

\begin{defn}[સુસંગતતા]
!!{sentence}sનો ગણ~$\Gamma$ \emph{અસુસંગત} છે જો અને માત્ર જો એવો સાન્ત
ગણ $\{!B_1, \dots, !B_n\} \subseteq \Gamma$ અને નીચેના ગણ માટે બંધ
!!{tableau} હોય:
\[
\{
  \sFmla{\True}{!B_1}, \dots,
  \sFmla{\True}{!B_n}
  \}.
\]
$\Gamma$ અસુસંગત ન હોય તો તેને \emph{સુસંગત} કહીએ છીએ.
\end{defn}

\begin{prop}[સ્વવાચકતા]
\ollabel{prop:reflexivity}
જો $!A \in \Gamma$ હોય, તો $\Gamma \Proves !A$.
\end{prop}

\begin{proof}
જો $!A \in \Gamma$ હોય, તો $\{!A\}$ એ~$\Gamma$નો સાન્ત ઉપગણ છે અને
નીચેનો !!{tableau}
\begin{oltableau}
  [\sFmla{\False}{\formula{A}}, just = \TAss
    [\sFmla{\True}{\formula{A}}, just = \TAss,close]
  ]
\end{oltableau}
બંધ છે.
\end{proof}


\begin{prop}[એકદિશવર્ધિતા]
\ollabel{prop:monotonicity}
જો $\Gamma \subseteq \Delta$ અને $\Gamma \Proves !A$ હોય, તો
$\Delta \Proves !A$.
\end{prop}

\begin{proof}
~$\Gamma$નો દરેક સાન્ત ઉપગણ~$\Delta$નો પણ સાન્ત ઉપગણ છે.
\end{proof}

\begin{prop}[પરંપરિતતા]
\ollabel{prop:transitivity}
જો $\Gamma \Proves !A$ અને $\{!A\} \cup \Delta \Proves !B$ હોય, તો
$\Gamma \cup \Delta \Proves !B$.
\end{prop}

\begin{proof}
જો $\{!A\} \cup \Delta \Proves !B$ હોય, તો એવો સાન્ત ઉપગણ
$\Delta_0 = \{!C_1, \dots, !C_n\} \subseteq \Delta$ છે કે
\begin{align*}
\{\sFmla{\False}{!B}, & \sFmla{\True}{!A}, \sFmla{\True}{!C_1},
\dots, \sFmla{\True}{!C_n}\}
\intertext{માટે બંધ !!{tableau} છે. જો $\Gamma \Proves !A$ હોય, તો એવો
  સાન્ત ઉપગણ $\{!D_1, \dots, !D_m\} \subseteq \Gamma$ છે કે}
\{\sFmla{\False}{!A}, & \sFmla{\True}{!D_1},
\dots, \sFmla{\True}{!D_m}\}
\end{align*}
માટે બંધ !!{tableau} છે.

હવે નીચેની ધારણાઓ ધરાવતો !!{tableau} લો:
\[
\sFmla{\False}{!B},
\sFmla{\True}{!C_1}, \dots, \sFmla{\True}{!C_n},
\sFmla{\True}{!D_1}, \dots, \sFmla{\True}{!D_m}.
\]
~$!A$ પર \Cut{} નિયમ લાગુ કરો. તેનાથી બે શાખાઓ બને છે: એકમાં
$\sFmla{\True}{!A}$ અને બીજીમાં $\sFmla{\False}{!A}$ હોય છે. તેથી એક
શાખામાં નીચેના બધા ઉપલબ્ધ છે:
\[
\{\sFmla{\False}{!B}, \sFmla{\True}{!A},
\sFmla{\True}{!C_1}, \dots, \sFmla{\True}{!C_n}\}
\]
આ ધારણાઓ માટે બંધ !!{tableau} હોવાથી તેને એ શાખા સાથે જોડી શકીએ છીએ;
$\sFmla{\True}{!A}$માંથી પસાર થતી દરેક શાખા બંધ થાય છે. બીજી શાખામાં
નીચેના બધા ઉપલબ્ધ છે:
\[
\{\sFmla{\False}{!A}, \sFmla{\True}{!D_1}, \dots,
\sFmla{\True}{!D_m}\}
\]
તેથી બીજી બાજુ પણ પૂરી કરીને બંધ !!{tableau} મેળવી શકીએ છીએ. આ
$\Gamma \cup \Delta \Proves !B$ દર્શાવે છે.
\sourcecorrection{OLTAB-003}{મૂળની \texttt{proof-theoretic-notions.tex},
પંક્તિ~105માં અલગ વાક્યો $!D_1,\dots,!D_m$ને સીધાં
``$\subseteq \Gamma$'' સાથે મૂક્યાં છે. સમાન વ્યાખ્યા અને તરત પછીના ગણ
મુજબ અહીં સાન્ત ગણ $\{!D_1,\dots,!D_m\}\subseteq\Gamma$ પુનઃસ્થાપિત
કર્યો છે.}
\end{proof}

નોંધો કે ખાસ કરીને તેનો અર્થ એવો થાય છે કે જો $\Gamma \Proves !A$ અને
$!A \Proves !B$ હોય, તો $\Gamma \Proves !B$. વળી, જો
$!A_1, \dots, !A_n \Proves !B$ અને દરેક~$i$ માટે
$\Gamma \Proves !A_i$ હોય, તો $\Gamma \Proves !B$ એવું પણ ફલિત થાય છે.

\begin{prop}
\ollabel{prop:incons}
$\Gamma$ અસુસંગત છે જો અને માત્ર જો દરેક !!{sentence}~$!A$ માટે
$\Gamma \Proves !A$.
\end{prop}

\begin{proof}
અભ્યાસપ્રશ્ન.
\end{proof}

\tagprob{FOL}
\begin{prob}
\olref[fol][tab][ptn]{prop:incons} સાબિત કરો.
\end{prob}
\tagendprob

\tagprob{notFOL}
\begin{prob}
\olref[pl][tab][ptn]{prop:incons} સાબિત કરો.
\end{prob}
\tagendprob

\begin{prop}[સઘનતા]
  \ollabel{prop:proves-compact}
  \begin{enumerate}
  \item જો $\Gamma \Proves !A$ હોય, તો એવો સાન્ત ઉપગણ
    $\Gamma_0 \subseteq \Gamma$ છે કે $\Gamma_0 \Proves !A$.
  \item જો $\Gamma$નો દરેક સાન્ત ઉપગણ સુસંગત હોય, તો $\Gamma$ સુસંગત છે.
  \end{enumerate}
\end{prop}

\begin{proof}
  \begin{enumerate}
    \item જો $\Gamma \Proves !A$ હોય, તો એવો સાન્ત ઉપગણ
      $\Gamma_0 = \{!B_1, \dots, !B_n\}$ અને નીચેના ગણ માટે બંધ
      !!{tableau} હોય છે:
      \[
      \{\sFmla{\False}{!A}, \sFmla{\True}{!B_1}, \dots, \sFmla{\True}{!B_n}\}
      \]
      આ !!{tableau} $\Gamma_0 \Proves !A$ પણ દર્શાવે છે.
    \item જો $\Gamma$ અસુસંગત હોય, તો કોઈ સાન્ત ઉપગણ
      $\Gamma_0 = \{!B_1, \dots, !B_n\}$ માટે નીચેના ગણનો બંધ
      !!{tableau} હોય છે:
      \[
      \{\sFmla{\True}{!B_1}, \dots, \sFmla{\True}{!B_n}\}
      \]
      આ બંધ !!{tableau} દર્શાવે છે કે $\Gamma_0$ અસુસંગત છે.
  \end{enumerate}
\end{proof}

\end{document}
